\section{Query Sequence Benchmark}
\label{sec:simulation-methodology}
In this section, we detail our simulation methodology and the resulting hyperparameters.
We base our sequence generator on SolidBench~\cite{taelman2023link}, 
which simulates realistic decentralized social media applications within the Solid ecosystem.
SolidBench utilizes the Social Network Benchmark (SNB) dataset~\cite{erling2015ldbc,angles2020ldbc} 
and applies the Linked Data Benchmark Council (LDBC) choke point methodology~\cite{angles2014linked}.

The choke point methodology is a benchmark design strategy that targets specific query execution 
and optimization bottlenecks in current database management systems~\cite{boncz2013tpc}.
Inspired by classical benchmarks like TPC-H, this approach uses these technical challenges to guide research and system development. 
For graph data, choke points include estimating cardinality during traversals with data skew, determining optimal join orders, 
and handling scattered index access patterns that lack predictable locality~\cite{erling2015ldbc}. 
In practice, a choke point-based benchmark design strategy first identifies critical execution challenges, 
then designs queries or modifications that explicitly test them.

We evaluate our prioritization algorithms using SolidBench~\cite{taelman2023link},
the state-of-the-art benchmark for structured decentralized environments.
SolidBench simulates a social network containing users, posts, and comments,
all described by a single ontology.
To represent Solid pods, the benchmark fragments data so that all information
related to a specific user resides in a single pod.
Under default settings, SolidBench generates 158,233 RDF files across
1,531 data vaults, totaling 3,556,159 triples.

To simulate a highly distributed network, SolidBench partitions the Social
Network Benchmark (SNB) dataset into individual data pods.
Its sequential architecture executes in four modular steps:
\begin{enumerate}
    \item The LDBC Datagen generates the underlying dataset.
    \item An enhancer creates additional data, such as noise triples that
    do not participate in queries.
    \item A dataset fragmenter partitions the triples into a Solid-like structure.
    \item The SPARQL queries are instantiated.
\end{enumerate}

The workload combines standard SNB \textit{interactive} queries (\textit{short} and \textit{complex})
with custom \textit{discover} templates that explicitly target
Link Traversal Query Processing (LTQP) choke points.
During execution, a client-side engine dynamically traverses links
between these scattered pods to resolve the queries.

To support sequence-based evaluation, we significantly extend and restructure three core SolidBench components:
\begin{itemize}
    \item \textbf{Query Generator:} Modified to produce correlated query sequences
    that reflect empirical usage patterns.
    \item \textbf{Dataset Generator:} Enhanced to model interest-based similarities
    among simulated users.
    \item \textbf{Benchmark Runner:} Updated with sequence execution logic to facilitate
    reproducible experiments and simplify the generation of default configuration files.
\end{itemize}

Link traversal engines cannot currently solve \textit{complex} queries in a feasible timeframe,
so we omit them from our default benchmark configuration.
However, because the sequence generator is fully configurable,
researchers can easily include complex query templates in future benchmarks
as traversal engines improve.


\subsection{Realism of Query Sequences}
To ensure SolidSessionBench accurately reflects user behavior in decentralized environments, 
we base our simulation methodology on the LDBC methodology used for data generation.
The sequence generator synthesizes empirical usage patterns established in Section~\ref{sec:related-work} to form a 
grounded simulation methodology that reflects real-world patterns.

Because direct, real-world query logs for highly decentralized LTQP environments do not currently exist,
we maintain and evaluate realism through a three-pillared approach:

\begin{itemize}
    \item \textbf{Topological Realism (User Interests):} We utilize the established LDBC social network topology
    to model explicit and implicit user relatedness, ensuring data access patterns mirror natural community formations 
    rather than uniform distributions. Due to the proven realism of the LDBC data generator~\cite{prat2014community}, 
    our simulations inherit this empirically proven structure.
    
    \item \textbf{Navigational Realism (Logical Sessions):} We restrict query sequences to directed dependency graphs.
    This simulates realistic ``click-through'' application navigation, 
    accounting for empirically observed session lengths and log-normal transition probabilities.  
 
    \item \textbf{Exploratory Realism (Query Refinement):} We replace static templates with dynamic query transformations. 
    Mapping real-world SPARQL refinements (e.g., result generalization, intent shifting) directly to execution choke points 
    ensures engines face realistic structural variability.
\end{itemize}

We evaluate the success of this simulation by measuring proxy variables such as
cache hit rates, execution times, and relatedness of relevant data accessed during queries. 
We then compare these outcomes against theoretical baselines and random sequence generation. 
This comparative evaluation demonstrates that our query sequence simulation captures essential 
performance dynamics that existing benchmarks miss.


\subsection{Simulation Methodology}
Our methodology models three interdependent behaviors: task-oriented logical sessions,
user interests for data locality,
and refinement patterns for iterative query development.

\subsubsection{Logical Query Sessions}
To model logical sessions, we categorize queries into four primary topics: 
person-specific messages, forum information, person attributes, and message attributes. 
Because queries in decentralized environments often depend on preceding results, 
we map the output types of each template to valid subsequent templates. 
This creates a directed graph of query progression that governs the flow of a logical session.
% \rt{If you can't add a figure in the introduction, I'd add one here, to illustrate what you explain in this paragraph.}

Within a session, a template's \emph{instantiation values} are typically derived 
from the previous query's results. 
To realistically simulate this ``click-through'' 
behavior, we execute centralized equivalents of the LDBC SNB queries using 
QLever~\cite{bast2017qlever}. 
We then randomly sample the returned result set 
to instantiate the subsequent query template. 
If a query yields no results, 
we fall back to our default, interest-based instantiation methodology.

As an example, consider the following data distributed across a 
decentralized environment:

\begin{lstlisting}[
    language=json,
    caption={(a) Alice's profile document (\url{https://pod.example/alice/profile.ttl}).},
    label={lst:example_dataset_alice_profile},
    breaklines=true,
    basicstyle=\small\ttfamily
]
<https://pod.example/alice/profile.ttl#me> a :Person ;
    :knows <https://pod.example/bob/profile.ttl#me> ;
    :hasInbox <https://pod.example/alice/inbox> .
\end{lstlisting}
\continuelisting

\begin{lstlisting}[
    language=json,
    caption={(b) Bob's profile document (\url{https://pod.example/bob/profile.ttl}).},
    label={lst:example_dataset_bob_profile},
    breaklines=true,
    basicstyle=\small\ttfamily
]
<https://pod.example/bob/profile.ttl#me> a :Person ;
    :posts <https://pod.example/bob/posts/1> .
\end{lstlisting}
\continuelisting

\begin{lstlisting}[
    language=json,
    caption={(c) Bob's post (\url{https://pod.example/bob/posts/1}).},
    label={lst:example_dataset_bob_post},
    breaklines=true,
    basicstyle=\small\ttfamily
]
<https://pod.example/bob/posts/1> a :Post ;
    :content "Hello Solid world!" .
\end{lstlisting}
\continuelisting

\begin{lstlisting}[
    language=json,
    caption={(d) Alice's interests document (\url{https://pod.example/alice/interests.ttl}).},
    label={lst:example_dataset_alice_interests},
    breaklines=true,
    basicstyle=\small\ttfamily
]
<https://pod.example/alice/profile.ttl#me> :interestedIn "Wildlife" .
\end{lstlisting}
\continuelisting

\begin{lstlisting}[
    language=json,
    caption={(e) Wildlife forum document (\url{https://pod.example/forum/wildlife.ttl}).},
    label={lst:example_dataset_forum_wildlife},
    breaklines=true,
    basicstyle=\small\ttfamily
]
<https://pod.example/forum/wildlife.ttl> a :Forum ;
    :hasPost <https://pod.example/ruben/posts/9> ;
    :hasComment <https://pod.example/bob/comment/1> ;
    :hasTopic "Wildlife" .
<https://pod.example/bob/comment/1> :responseTo <https://pod.example/ruben/posts/9> .
\end{lstlisting}

% \begin{lstlisting}[
%     language=json,
%     caption={A minimal decentralized dataset distributed across documents.},
%     label={lst:example_dataset},
%     breaklines=true,
%     basicstyle=\small\ttfamily
% ]
% @prefix : <https://example.org/voc#> .

% # File: https://pod.example/alice/profile.ttl
% <https://pod.example/alice/profile.ttl#me> a :Person ;
%     :knows <https://pod.example/bob/profile.ttl#me> ;
%     :hasInbox <https://pod.example/alice/inbox> .

% # File: https://pod.example/bob/profile.ttl
% <https://pod.example/bob/profile.ttl#me> a :Person ;
%     :posts <https://pod.example/bob/posts/1> .

% # File: https://pod.example/bob/posts/1
% <https://pod.example/bob/posts/1> a :Post ;
%     :content "Hello Solid world!" .

% # File: https://pod.example/alice/interests.ttl
% <https://pod.example/alice/profile.ttl#me> :interestedIn "Wildlife" .

% # File: https://pod.example/forum/wildlife.ttl
% <https://pod.example/forum/wildlife.ttl> a :Forum ;
%     :hasPost <https://pod.example/ruben/posts/9> ;
%     :hasComment <https://pod.example/bob/comment/1> ;
%     :hasTopic "Wildlife" .
% <https://pod.example/bob/comment/1> :responseTo <https://pod.example/ruben/posts/9> .
% \end{lstlisting}

The generator randomly selects a user and simulates
a realistic sequence of queries that reflects real-world user behavior.
Given 
Given this dataset, and two query 
templates: one to fetch a person's posts and another to retrieve the 
content of a specific post, the generator can create a query sequence, which can contain multiple
logical sessions.
For user ``alice'', we can start the sequence (and first logical session) by 
executing the first query template (Q1) to discover Bob's posts, 
as shown in Listing~\ref{lst:query_1}.

\begin{lstlisting}[
    language=SPARQL,
    caption={Initial query (Q1) fetching Bob's posts.},
    label={lst:query_1},
    breaklines=true,
    basicstyle=\small\ttfamily
]
SELECT ?post WHERE {
  <https://pod.example/bob/profile.ttl#me> :posts ?post .
}
\end{lstlisting}

The client-side query engine executes this query by traversing to Bob's 
profile document, returning the result binding: 
`?post = <https://pod.example/bob/posts/1>`.

To simulate a user clicking on this discovered post, the sequence 
generator instantiates the next query template dynamically. It extracts 
the value bound to ``?post'' from the results of Q1 and injects it as the 
subject of Q2, as shown in Listing~\ref{lst:query_2}.

\begin{lstlisting}[
    language=SPARQL,
    caption={Dynamic sequence query (Q2) instantiating a variable from Q1.},
    label={lst:query_2},
    breaklines=true,
    basicstyle=\small\ttfamily
]
SELECT ?content WHERE {
  <https://pod.example/bob/posts/1> :content ?content .
}
\end{lstlisting}

This example demonstrates that the instantiation of Q2 strictly depends 
on the execution results of Q1. Because Q2 requires an instantiated 
post URI from Q1, these queries are directionally dependent; reversing 
their order leaves the generator with no method to instantiate the 
initial query parameters. Consequently, structured logical query sequences 
naturally emerge from these interdependencies, accurately reflecting 
realistic exploratory user behavior.

Finally, we simulate logical session switching. 
Drawing from web browsing behavior, we assign each sequence a session-switching probability sampled from a log-normal 
distribution with a mean of 15\%~\cite{meiss_whats_2009,schneider2009understanding}. 
During sequence generation, a session switch triggers a random choice with equal probability: the engine either starts a new 
session (using interest-based instantiation) or resumes a previous one (deriving instantiation values from that session's last executed query).
To prevent repetitive queries and maintain a uniform distribution of template instantiations, 
we explicitly prohibit backtracking within our model. 
Additionally, we maintain a balanced template distribution by dynamically 
scaling down each template's selection probability proportionally to its 
overall frequency.

Continuing our running example, a session switch occurs after executing Q2. 
Instead of continuing sequence generation from the content of Bob's post, 
the generator terminates the current path and selects a new template to 
instantiate: a query for posts within a forum.

As shown in Listing~\ref{lst:query_3}, the generator selects a wildlife forum 
as the initial instantiation value. This query requires an entirely separate 
data path to execute, breaking the dependency chain of the previous queries 
and representing a clear change in user intent. We are still generating the same
\textit{query sequence}, but the session switch introduces a new \textit{logical session} within that sequence.

\begin{lstlisting}[
    language=json,
    caption={Independent query (Q3) initiating a session switch to explore a wildlife forum.},
    label={lst:query_3},
    breaklines=true,
    basicstyle=\small\ttfamily
]
SELECT ?post WHERE {
  <https://pod.example/forum/wildlife.ttl> :hasPost ?post .
}
\end{lstlisting}

\subsubsection{User Interest Modeling}

Individuals interact more frequently with users sharing similar interests~\cite{mitzlaff2013user}. 
To simulate this behavior, we use the underlying LDBC dataset utilized by SolidBench~\cite{erling2015ldbc}, 
which inherently models realistic user preferences. 
The LDBC generator, Datagen, produces person profiles containing specific interests, 
educational backgrounds, and workplaces. 
It links users based on these demographic and interest overlaps, 
yielding a non-random network topology with a realistic degree distribution.

We use this interest-driven topology to generate template instantiation values 
for the first query of each logical session. 
Specifically, we construct a binary feature vector $v_{k}$ for each user $k$:
\begin{equation}
 v_{k} = (x_{k,1}, \dots, x_{k,n}, y_{k,1}, \dots, y_{k,m}),
\end{equation}
where $x_{k,j}=1$ if user $k$ holds stated interest $j$, and $y_{k,m}=1$ 
if user $k$ completed study $m$.

For each sequence, we randomly select a base person $k$ and compute the cosine similarity 
between $v_{k}$ and the feature vectors of all other users.
Let $N$ be the total number of users. To maintain efficiency at large scale factors, 
we restrict our candidate pool to $C_{k,M}$, defined as the set of the top $M < N$ users 
ranked by similarity to $k$. 
Template instantiation values are then selected probabilistically using a softmax function 
with temperature $T$:

\begin{equation}
   P(i \mid k) = \frac{e^{s_{k,i} / T}}{\sum_{u \in C_{k,M}} e^{s_{k,u} / T}},
\end{equation}

where $s_{k,i}$ is the similarity score between base person $k$ and candidate $i \in C_{k,M}$. 
For templates requiring non-person values (e.g., posts or activities),
selection probabilities are derived from the similarity scores of the content's creators.

Continuing our running example, consider an expanded decentralized environment containing a
second forum dedicated to pottery, as shown in
Listing~\ref{lst:augmented_dataset} as well as all documents in Listing~\ref{lst:example_dataset_alice_interests}.

\begin{lstlisting}[
    language=json,
    caption={The additional pottery forum document (\url{https://pod.example/forum/pottery.ttl}).},
    label={lst:augmented_dataset},
    breaklines=true,
    basicstyle=\small\ttfamily
]
<https://pod.example/forum/pottery.ttl> a :Forum ;
    :hasTopic "Pottery" .
\end{lstlisting}

To determine what instantiation value to use for the forum query for user Alice, we construct 
binary feature vectors over a small global feature space: 
$\text{(Wildlife, Pottery)}$. 

Alice's profile states an explicit interest in `Wildlife', yielding vector 
$v_{\text{Alice}} = (1, 0)$. While the forums in the data have topic statements, thus
the vectors for the forums are $v_{\text{Wildlife}} = (1, 0)$ and
$v_{\text{Pottery}} = (0, 1)$. 

This gives us a cosine similarity of $s_{\text{Alice},\text{Wildlife}} = 1$ and
$s_{\text{Alice},\text{Pottery}} = 0$. Then using softmax with a temperature of $T = 0.5$,
we compute the selection probabilities for the two forums:

\begin{align*}
    P(\text{Wildlife} \mid \text{Alice}) = \frac{e^{1 / 0.5}}{e^{1 / 0.5} + e^{0 / 0.5}} \approx 0.88, \\
    P(\text{Pottery} \mid \text{Alice}) = \frac{e^{0 / 0.5}}{e^{1 / 0.5} + e^{0 / 0.5}} \approx 0.12.
\end{align*}

Consequently, when initiating or switching a session, the generator is more likely to select the
`Wildlife' forum for Alice, reflecting her interests.

\subsubsection{Refinement Patterns}

The benchmark simulates query refinement by applying 
composable transformations to an instantiated query template. 
Following our literature review, we focus on transformations 
(and their reciprocal counterparts) that add or modify 
BGPs, \texttt{OPTIONAL} clauses, \texttt{UNION} blocks, 
and query instantiation values.

To ensure these transformations meaningfully alter query execution 
rather than superficially appending triples, we apply a 
choke point-based design methodology~\cite{angles2014linked}. 
The generator sequentially applies refinements to a base 
SolidSessionBench template. By mapping empirical exploratory 
usage patterns~\cite{arenas2025exploring} to concrete 
planning (P) and traversal (T) choke points, we translate 
real-world user behaviors into rigorous execution bottlenecks 
for query engines.

The planning choke points, derived from empirical usage 
patterns~\cite{arenas2025exploring} and relevant LDBC 
choke points~\cite{angles2020ldbc}, are defined as follows:
\begin{enumerate}[label=\textbf{P.\arabic*}, leftmargin=*]
    \item Involves adding or 
    removing a \texttt{FILTER}, requiring prior knowledge 
    to adjust query plans and estimate selectivity.
    \item Addresses 
    adding or removing \texttt{UNION} or \texttt{OPTIONAL} 
    clauses to alter filter push-down capabilities and 
    complicate query planning.
    \item Covers adding 
    or removing selective triple patterns, forcing the engine 
    to identify changes and re-optimize the plan.
    \item Focuses 
    on adding or removing non-selective triple patterns.
\end{enumerate}

Beyond planning, these refinements introduce traversal 
choke points that expand upon those established in 
SolidBench~\cite{taelman2023link,taelman2023evaluation}:
\begin{enumerate}[label=\textbf{T.\arabic*}, leftmargin=*]
    \item Expands 
    or contracts the reachable sub-web through 
    predicate additions.
    \item Shifts the sub-web 
    by substituting the traversal starting point.
    \item Changes the number 
    of hops required to obtain relevant data, corresponding 
    to SolidBench choke points L.1--L.6~\cite{taelman2023evaluation}.
    \item Alters the usability of structural indexes, such as type indexes, 
    corresponding to SolidBench choke point 
    L.8~\cite{taelman2023evaluation}.
\end{enumerate}

These choke points directly map to the exploratory query 
behaviors identified in Section~\ref{sec:related_work_qup}:
\begin{itemize}[leftmargin=*]
    \item \textbf{Result refinement} corresponds to 
    P.1 and P.3, as it narrows results via \texttt{FILTER} 
    constraints and selective triple patterns.
    \item \textbf{Result generalization} corresponds to 
    P.3 and P.4, expanding existing BGPs with selective 
    or non-selective triple patterns.
    \item \textbf{Result extension} maps to P.2, P.3, 
    and P.4 by appending triple patterns to BGPs or 
    introducing \texttt{OPTIONAL} and \texttt{UNION} 
    blocks that alter filter push-down capabilities.
    \item \textbf{Shifting query intent} corresponds to 
    T.1 and T.2, substituting instantiation values to 
    initiate traversal from different seed documents, 
    thereby shifting or expanding the reachable sub-web.
    \item \textbf{Query refinement} maps to P.2, 
    modifying the structural complexity of the query 
    by adding or removing \texttt{OPTIONAL} and 
    \texttt{UNION} blocks.
\end{itemize}

Finally, the generator supports variable instantiation 
within refinements, allowing users to bind pattern variables 
to specific values randomly selected from configuration candidates. 
Users can define custom refinement patterns via JSON files to 
augment the default configurations. 

An example of a slimmed-down default configuration is 
provided in Listing~\ref{lst:refinement_config}. 
This entry specifies an \texttt{OPTIONAL} refinement 
pattern that appends a triple pattern to retrieve the 
messages liked by a person. 

The \texttt{location} parameter is set to 0, indicating 
that the triple pattern is added to the first existing 
\texttt{OPTIONAL} block, or to a new block if none is 
present. The default benchmark suite also includes a 
reciprocal configuration where \texttt{operation} is 
set to ``removal'', which instead deletes the triple pattern 
from the first \texttt{OPTIONAL} block.

\begin{lstlisting}[
    language=json,
    caption={Example configuration for an \texttt{OPTIONAL} refinement pattern.},
    label={lst:refinement_config},
    breaklines=true,
    basicstyle=\small\ttfamily,
    columns=fullflexible
]
{
  "type": "OPTIONAL",
  "id": 0,
  "operation": "addition",
  "description": "Add triple for likes (P.1, T.2)",
  "location": 0,
  "target": [
    {
      "subject": { 
        "value": "person", 
        "termType": "variable" 
      },
      "predicate": { 
        "value": ".../vocabulary/likes", 
        "termType": "namedNode" 
      },
      "object": { 
        "value": "likedMessage", 
        "termType": "variable" 
      }
    }
  ]
}
\end{lstlisting}

All defaults, along with 
complete documentation, are available on GitHub.%
\footnote{\url{https://github.com/SolidBench/SolidBench.js/tree/master/templates/refinements}}


\subsection{Hyperparameters}

The sequence generation process introduces several new hyperparameters, 
which are detailed in Table~\ref{tab:generation_hyperparameters} together with their default values. 
To guarantee reproducibility despite extensive random sampling, 
a single seeded pseudo-random number generator (PRNG) is used througout the 
entire pipeline.

\begin{table}[htbp]
    \centering
    \caption{LogT-normal distribution parameters ($\mu, \sigma$) control the queries per sequence, 
    queries per logical session, and the probability of switching logical sessions. 
    Refinement parameters define the likelihood (\texttt{patternProbability}) 
    and bounds (\texttt{min/maxRefinementLength}) of generating a refinement sequence for a given query. 
    Instantiation relies on a softmax \texttt{temperature} and a \texttt{maxLogits} 
    cutoff for similarity-based entity selection.}
    \label{tab:generation_hyperparameters}
    \begin{tabular}{l l c}
        \toprule
        \textbf{Category} & \textbf{Parameter} & \textbf{Value} \\
        \midrule
        \textbf{Distributions} ($\mu, \sigma$) 
        & \texttt{Sequence Length} & 3.0, 0.2 \\
        & \texttt{Session Length} & 1.7, 0.5 \\
        & \texttt{Transition Probability} & -2.0, 0.25 \\
        \addlinespace
        \textbf{Refinements} 
        & \texttt{patternProbability} & 0.1 \\
        & \texttt{min/maxRefinementLength} & 2, [3--5] \\
        \addlinespace
        \textbf{Instantiation} 
        & \texttt{temperature} & 0.1 \\
        & \texttt{maxLogits} & 5--10 \\
        \bottomrule
    \end{tabular}
\end{table}

